\section*{Conservation vector}
Set
\[
c=(0,\underbrace{4,\ldots,4}_{X_2,\ldots,X_{m-1}},2,1)^T.
\]
For every reaction $y\to y'$ in $\widehat N_m$, direct inspection gives
$c^T(y'-y)=0$: the inflow creates $X_1$, which has weight zero; each chain
step replaces one weight-four species by another; the reaction
$X_1+X_{m-1}\to2X_m$ preserves $4=2+2$; $2X_m\to X_2$ preserves $4$; and
$2Z\rightleftarrows X_1+X_m$ preserves $2=0+2$.  Hence $c^T\Gamma_m=0$.


\section*{Stoichiometric space}
The $m$ long-circuit columns contain one dependence with all coefficients one;
the reversible pair contributes a second, independent dependence.  Exact row
reduction gives $\operatorname{rank}\Gamma_m=m$.  Since $c\neq0$ is a left
null vector in $\mathbb R^{m+1}$, the image has codimension one and therefore
\[
\operatorname{im}\Gamma_m=c^\perp.
\]
At the selected realization the right kernel of $A_m$ is generated by
\[
\rho=(2,-2,\ldots,-2,0,1)^T,
\qquad c^T\rho=17-8m\neq0.
\]
Consequently $\ker A_m\cap c^\perp=\{0\}$, so the homogeneous restriction to
the stoichiometric space has no conservation zero.


\section*{Fixed integrated mass phase space}
On $(0,\pi)$ with Neumann boundary conditions define
\[
X_c=\left\{u\in H_N^2((0,\pi);\mathbb R^{m+1}):
\int_0^\pi c^Tu\,d\xi=0\right\},
\]
\[
Y_c=\left\{g\in L^2((0,\pi);\mathbb R^{m+1}):
\int_0^\pi c^Tg\,d\xi=0\right\}.
\]
Because $c^Tf(x)=0$ pointwise and
$\int_0^\pi c^TDu''=c^TDu'|_0^\pi=0$, the stationary operator maps
$X_c$ into $Y_c$.  The critical function $r\cos\xi$ lies in $X_c$ even when
$c^Tr\neq0$, because $\cos\xi$ has zero mean.  Only the homogeneous Fourier
coefficient is constrained to lie in $c^\perp$.
